% =============================================================================
% style/base_section.tex
% =============================================================================


% =============================================================================
% Цвета заголовков разделов
% =============================================================================

%\colorlet{sectionAccent}{xindigo!78!xblue}
\colorlet{sectionAccent}{xindigo!55!black}
\colorlet{sectionAccentDark}{sectionAccent!55!black}

% Цвета номерного бейджа.
\colorlet{sectionBadgeTop}{sectionAccent!72!black}
\colorlet{sectionBadgeBottom}{sectionAccent!28!black}
\colorlet{sectionBadgeFrame}{sectionAccent!38!black}
\colorlet{sectionBadgeText}{sectionAccent!6!white}

% Цвет заголовка и двух нижних линий.
\colorlet{sectionTitleColor}{sectionAccentDark}
\colorlet{sectionRuleA}{sectionAccent!65!white}
\colorlet{sectionRuleB}{sectionAccent!45!white}

% Две вертикальные линии для ненумерованного section.
\colorlet{sectionNumberlessLineA}{sectionAccent!65!white}
\colorlet{sectionNumberlessLineB}{sectionAccent!75!black}


% =============================================================================
% Fading для чистого верхнего блика
% =============================================================================

\tikzfading[
  name=sectionGlossFade,
  top color=transparent!0,
  bottom color=transparent!100,
]


% =============================================================================
% Геометрия
% =============================================================================

\newlength{\sectionBadgeBaselineShift}
\setlength{\sectionBadgeBaselineShift}{5.2pt}

\newlength{\sectionNumberlessBaselineShift}
\setlength{\sectionNumberlessBaselineShift}{5.2pt}


% =============================================================================
% Номерной глянцевый бейдж для section
% =============================================================================

\providecommand{\sectionGlossyBadge}[1]{}
\renewcommand{\sectionGlossyBadge}[1]{%
  \tikz[baseline={([yshift=-\sectionBadgeBaselineShift]S.center)}]{%
    % Невидимый узел задаёт размеры бейджа.
    \node[
      minimum width=13.6mm,
      minimum height=7.8mm,
      inner xsep=4.4pt,
      inner ysep=1.3pt,
      outer sep=0pt,
      font=\sffamily\bfseries\fontsize{10.8}{12}\selectfont
    ] (S) {\phantom{\strut #1}};

    % Внешняя тень.
    \begin{scope}[xshift=0.38mm, yshift=-0.46mm]
      \path[
        fill=black,
        opacity=0.16,
        rounded corners=2.0pt,
      ]
        (S.south west) rectangle (S.north east);
    \end{scope}

    % Основная заливка: чистый вертикальный градиент.
    \shade[
      top color=sectionBadgeTop,
      bottom color=sectionBadgeBottom,
      draw=sectionBadgeFrame,
      line width=0.58pt,
      rounded corners=2.0pt,
    ]
      (S.south west) rectangle (S.north east);

    % Верхний глянцевый блик: плавно исчезает к середине бейджа.
    \begin{scope}
      \clip[rounded corners=2.0pt]
        (S.south west) rectangle (S.north east);
      \path[
        fill=white,
        path fading=sectionGlossFade,
        opacity=0.52,
      ]
        ([yshift=-0.18mm]S.north west) rectangle
        ([yshift=-0.18mm]S.east);
    \end{scope}

    % Верхний внутренний световой кант.
    \begin{scope}
      \clip[rounded corners=2.0pt]
        (S.south west) rectangle (S.north east);
      \draw[
        white,
        opacity=0.55,
        line width=0.40pt,
        line cap=round,
      ]
        ([xshift=1.1mm, yshift=-0.38mm]S.north west) --
        ([xshift=-1.1mm, yshift=-0.38mm]S.north east);
    \end{scope}

    % Нижний внутренний теневой кант.
    \begin{scope}
      \clip[rounded corners=2.0pt]
        (S.south west) rectangle (S.north east);
      \draw[
        black,
        opacity=0.20,
        line width=0.38pt,
        line cap=round,
      ]
        ([xshift=1.1mm, yshift=0.40mm]S.south west) --
        ([xshift=-1.1mm, yshift=0.40mm]S.south east);
    \end{scope}

    % Обводка поверх бликов.
    \path[
      draw=sectionBadgeFrame,
      line width=0.58pt,
      rounded corners=2.0pt,
    ]
      (S.south west) rectangle (S.north east);

    % Номер section.
    \node[
      minimum width=13.6mm,
      minimum height=7.8mm,
      inner xsep=4.4pt,
      inner ysep=1.3pt,
      outer sep=0pt,
      text=sectionBadgeText,
      font=\sffamily\bfseries\fontsize{10.8}{12}\selectfont,
    ] at (S.center) {\strut #1};
  }%
}


% =============================================================================
% Метка ненумерованного section: две вертикальные линии
% =============================================================================

\providecommand{\sectionNumberlessMark}{}
\renewcommand{\sectionNumberlessMark}{%
  \tikz[baseline={([yshift=-\sectionNumberlessBaselineShift]S.center)}]{%
    \node[
      minimum width=5.0mm,
      minimum height=7.8mm,
      inner sep=0pt,
      outer sep=0pt
    ] (S) {};

    \draw[
      sectionNumberlessLineA,
      line width=1.30pt,
      line cap=round,
    ]
      ([xshift=1.4mm, yshift=-0.2mm]S.north west) --
      ([xshift=1.4mm, yshift=0.2mm]S.south west);

    \draw[
      sectionNumberlessLineB,
      line width=1.50pt,
      line cap=round,
    ]
      ([xshift=3.1mm, yshift=-0.2mm]S.north west) --
      ([xshift=3.1mm, yshift=0.2mm]S.south west);
  }%
}


% =============================================================================
% Две линии под нумерованным section
% =============================================================================

\providecommand{\sectionDoubleRule}{}
\renewcommand{\sectionDoubleRule}{%
  \vspace{3.0pt}%
  {\color{sectionRuleA}\titlerule[1.00pt]}%
  \vspace{1.25pt}%
  {\color{sectionRuleB}\titlerule[0.65pt]}%
}


% =============================================================================
% Нумерованные section: бейдж + две горизонтальные линии
% =============================================================================

\titleformat{\section}[hang]
  {%
    \normalfont\sffamily\bfseries\raggedright\color{sectionTitleColor}%
    \hyphenpenalty=10000\exhyphenpenalty=10000%
  }
  {\sectionGlossyBadge{\thesection}}
  {10pt}
  {\fontsize{18}{22}\selectfont}
  [\sectionDoubleRule]


% =============================================================================
% Ненумерованные section: только вертикальная метка, без горизонтальных линий
% =============================================================================

\titleformat{name=\section,numberless}[hang]
  {%
    \normalfont\sffamily\bfseries\raggedright\color{sectionTitleColor}%
    \hyphenpenalty=10000\exhyphenpenalty=10000%
  }
  {\sectionNumberlessMark}
  {10pt}
  {\fontsize{18}{22}\selectfont}


% =============================================================================
% Отступы section
% =============================================================================

\titlespacing*{\section}
  {0pt}
  {5.7ex plus 1.8ex minus 1.4ex}
  {2.6ex plus 0.6ex minus 0.5ex}